\chapter*{Introduction}
\addcontentsline{toc}{chapter}{Introduction}

\begin{longquote}
    Les Buts de la Documentation organisée consistent à pouvoir offrir sur tout ordre de fait et de connaissance des informations documentées : 1\textdegree{} universelles quant à leur objet ; 2\textdegree{} sûres et vraies ; 3\textdegree{} complètes ; 4\textdegree{} rapides ; 5\textdegree{} à jour ; 6\textdegree{} faciles à obtenir ; 7\textdegree{} réunies d'avance et prêtes à être communiquées ; 8\textdegree{} mises à la disposition du plus grand nombre. \parencite{otletTraiteDocumentationLivre2021}
\end{longquote}

\begin{longquote}
    Les archives constituent un patrimoine unique et irremplaçable transmis de génération en génération [...].
    Sources d'informations fiables pour une gouvernance responsable et transparente, [elles] jouent un rôle essentiel dans le développement des sociétés. \parencite{UniversalDeclarationArchives}
\end{longquote}

Ces deux textes définissent le contrat social de l'archivistique moderne : transformer la trace brute en connaissance organisée, universellement accessible et durablement fiable.
Pres de un siècle après Otlet, cette ambition s'est déplacée vers le Web sémantique et les graphes de connaissances --- seuls capables aujourd'hui de tenir la promesse d'une information structurée au-delà du document isolé.
En Afrique subsaharienne pourtant, cet idéal se heurte à une réalité matérielle qui impose d'intégrer une dimension absente des paradigmes dominants des pays Nord : la frugalité.

\subsubsection{Contexte et motivation}

L'Afrique subsaharienne connaît une accélération de la transformation numérique de ses administrations publiques \parencite{lawalDIGITALGOVERNANCEINCLUSIVE2025,egbeleoDigitalTransformationInstitutional2025}.
Dans cette dynamique, l'enregistrement des faits d'état civil s'impose comme une infrastructure de gouvernance, d'inclusion sociale et de souveraineté nationale \parencite{BirthRegistrationSubSaharan2025} : garantir à chaque individu une existence légale, conformément à l'Objectif de Développement Durable 16.9 des Nations Unies, conditionne toute planification en matière d'éducation, de santé et de droits fondamentaux.
Or la gestion documentaire de la région fait face à une crise silencieuse d'amnésie de ses passifs archivistiques : la production administrative, bien qu'indispensable, a longtemps été gérée de manière inadéquate, mettant à risque la préservation de la mémoire collective et la justification des droits individuels \parencite{michaudSecretsDEtats2015, bintsenaArchivesAdministrativesDans2017}.

Plusieurs pays de la sous-région illustrent cette tension entre ambition de modernisation et réalité de terrain.
Le Sénégal est parvenu à numériser plus de 20 millions d'actes via le projet NEKKAL \parencite{civipol2023nekkal}, ce modèle reposait sur des moyens financiers et humains extrêmement lourds --- subventions internationales et indexation manuelle massive ---, rendant sa reproduction problématique pour les États aux ressources plus limitées.
Le Burkina Faso a déployé iCivil, un système d'enregistrement nativement numérique --- tourné vers le flux des actes à venir, non vers l'indexation rétrospective du stock hérité.
En 2024, le Cameroun a engagé un tournant normatif majeur par l'adoption simultanée de lois régissant les archives numériques (n° 2024/001), l'état civil (n° 2024/016) et la protection des données personnelles (n° 2024/017).
Cependant, ce nouvel élan juridique se heurte à d'importantes contraintes du terrain : un déficit documentaire touchant près de 7 millions d'individus et la sous-exploitation d'un patrimoine d'actes manuscrits anciens restreint à des \og îlots informationnels \fg{} locaux \parencite{olembe2024webinaire}.
Nous n'avons pas identifié, dans la littérature et les sources institutionnelles consultées, de dispositif documenté combinant aujourd'hui numérisation de masse et indexation sémantique du stock manuscrit hérité --- c'est précisément ce vide que ce mémoire investit.

Cette même tension se lit au niveau des standards de description archivistique. Les modèles hérités --- OAIS, EAD, Dublin Core --- structurent un fonds selon une logique arborescente à racine unique, inapte à représenter un acte d'état civil mettant simultanément en jeu plusieurs agents et référencé par d'autres actes ultérieurs.
Le standard Records in Contexts (RiC) publié par le Conseil international des archives, dépasse cette rigidité en substituant à l'arbre un graphe de relations typées --- mais cette bascule ne résout ni l'indexation du contenu de ces documents manuscrits, ni le coût des ressources qu'exige son déploiement effectif.

Les systèmes documentaires classiques reposent majoritairement sur une indexation lexicale ou contrôlée, fondée sur des mots-clés, des descripteurs, des classifications ou des thésaurus \parencite{menilletThesaurusIndexation1993}, ce qui les rend moins aptes à représenter des relations complexes se limitant à retrouver des documents contenant certains mots, non des faits : une requête portant sur une même lignée familiale n'y a simplement pas de formulation possible.
L'indexation sémantique par graphe de connaissances offre une réponse structurelle en représentant chaque personne comme un nœud persistant à travers les actes.
Cette expressivité a un double coût : un graphe ne dissout pas le bruit de la reconnaissance manuscrite, il le déplace vers un problème de résolution d'entités \parencite{alkendiAdvancementsChallengesHandwritten2024,tuselmannNamedEntityLinking2022}, et son déploiement mobilise une infrastructure de calcul dont la plupart des services d'archives subsahariens ne disposent pas, au risque d'une dépendance technologique exogène menaçant la souveraineté numérique des États sur leurs données d'identité \parencite{musoniDigitalIDSystems}.


\subsubsection{Question de recherche principale}
C'est de cette confrontation --- une solution conceptuellement adéquate mais de ressource inaccessible --- que ce mémoire pose ainsi la question suivante :

\emph{Dans quelle mesure l'intégration de techniques de traitement automatique des langues frugales permet-elle de valider la viabilité d'un modèle de graphe de connaissances normalisé (RiC-O) pour la préservation et l'exploitation des archives manuscrites d'état civil en Afrique subsaharienne, en milieu à ressources computationnelles limitées ?}

Plus spécifiquement, ce mémoire évalue si une mobilisation des modèles de langue frugaux plutôt que leurs homologues industriels, des méthodes d'adaptation elles-mêmes économes en données annotées et une exécution intégrale sur infrastructure strictement frugale est compatible avec une qualité sémantique satisfaisante.

\subsubsection{Questions de recherche secondaires}

Cette viabilité se décline en quatre Questions de recherche (QRs) :



: quelle est la corrélation entre la variabilité des écritures cursives administratives subsahariennes et la fiabilité des transcriptions automatiques au sein d'un pipeline frugal ? (Q1\label{ques:1}). La deuxième interroge la modélisation : comment adapter l'ontologie RiC-O aux structures de parenté et de filiation propres aux contextes subsahariens sans rompre l'interopérabilité internationale du standard ? (Q2\label{ques:2}). La troisième s'attaque au cœur du compromis technique : quelles sont les limites de performance des modèles de langue frugaux dans l'extraction de relations complexes depuis un texte administratif manuscrit bruité ? (Q3\label{ques:3}). Enfin, la dernière évalue une stratégie de mitigation : une architecture d'extraction hybride, combinant règles symboliques et modèle statistique frugal, améliore-t-elle la robustesse du pipeline par rapport à un modèle neuronal utilisé seul ? (Q4\label{ques:4}).

À ces questions répondent quatre hypothèses de recherche, conçues comme des postulats falsifiables. Un graphe de connaissances conforme à RiC-O réduit significativement les silences documentaires liés aux variantes onomastiques, comparativement à une indexation lexicale classique (H1\label{hypo:1}). Une quantification INT8 des modèles frugaux mobilisés n'entraîne pas de dégradation critique de leur performance d'extraction, tout en rendant possible une exécution intégrale sur infrastructure CPU (H2\label{hypo:2}). Un profil d'application RiC-O simplifié et adapté aux structures de parenté subsahariennes représente des fonds hétérogènes sans recourir à une centralisation matérielle lourde (H3\label{hypo:3}). Enfin, une stratégie hybride combinant règles symboliques et modèle statistique frugal améliore la robustesse de l'extraction sur texte administratif bruité, comparativement à un modèle neuronal frugal utilisé seul (H4\label{hypo:4}).

Conformément au paradigme de la Design Science Research \parencite{hevnerDesignScienceInformation,Peffers2007ADS}, cette recherche distingue des objectifs scientifiques, orientés vers la production de connaissances généralisables, et des objectifs opérationnels, orientés vers la conception de l'artefact qui sert de support empirique à cette évaluation. Les objectifs scientifiques sont au nombre de trois : caractériser les limites de scalabilité sémantique de RiC-O dans un environnement de calcul dégradé appliqué au patrimoine manuscrit subsaharien (OS1) ; modéliser un profil RiC-O adaptable aux particularités juridiques, onomastiques et parentales des administrations d'Afrique subsaharienne (OS2) ; et évaluer empiriquement l'impact du bruit de transcription HTR sur la robustesse des requêtes SPARQL et sur la qualité du graphe produit (OS3). Les objectifs opérationnels visent la conception de l'artefact : concevoir et implémenter, sur un corpus de preuve de concept de 50 à 100 actes manuscrits camerounais numérisés par photographie, un pipeline frugal complet articulant HTR, extraction d'entités et de relations, et peuplement d'un graphe RiC-O (OO1) ; fine-tuner et quantifier des modèles frugaux exécutables sur infrastructure strictement CPU, interrogeables via un index vectoriel comparé à un moteur lexical de référence (OO2) ; et intégrer, en amont du peuplement du graphe, un mécanisme de pseudonymisation conforme à la loi camerounaise n° 2024/017, comme exigence de conception plutôt que comme contrainte ajoutée après coup (OO3).

L'originalité de cette recherche réside dans la confrontation expérimentale inédite du standard RiC \parencite{clavaudICARecordsContextsOntologyb} avec un environnement de traitement par intelligence artificielle local, souverain et frugal, appliqué au patrimoine manuscrit d'état civil subsaharien. Cette intersection n'est couverte ni par les études de numérisation archivistique africaine --- centrées sur la préservation plutôt que sur l'indexation sémantique ---, ni par les architectures sémantiques appliquées à l'état civil ailleurs dans le monde --- développées sans contrainte de calcul frugale ---, ni par les innovations numériques africaines en matière d'enregistrement --- tournées vers le flux à venir plutôt que vers le stock hérité. Cette ambition s'accompagne cependant de limites assumées. L'étude porte sur un corpus restreint de langue française, à l'échelle d'une preuve de concept de 50 à 100 actes camerounais, sans ambition de validation industrielle : elle privilégie la démonstration de viabilité d'un modèle conceptuel. Le profil RiC-O simplifié proposé ne prétend pas non plus se substituer à une adaptation exhaustive du standard à la diversité des pratiques de parenté d'Afrique subsaharienne. Par ailleurs, le caractère multidisciplinaire de ce travail --- au croisement de l'archivistique, du traitement automatique des langues et de l'informatique frugale --- comporte le risque de ne pas approfondir autant chaque domaine séparé, pour privilégier une intégration globale dont l'efficacité sera validée par la mise en pratique du prototype.

Le présent mémoire s'organise en quatre chapitres, qui suivent la trajectoire de la Design Science Research annoncée ci-dessus. Le premier chapitre situe la recherche dans la littérature et identifie les lacunes qu'elle comble. Le deuxième construit le dispositif méthodologique : positionnement épistémologique, profil RiC-O, architecture du pipeline, stratégie de pseudonymisation et protocole expérimental. Le troisième rend compte de l'implémentation et des résultats empiriques obtenus à chaque étage de la chaîne. Enfin, le quatrième interprète ces résultats à la lumière de la problématique, discute les limites du dispositif et ouvre sur les perspectives de déploiement à l'échelle sous-régionale.